\section{Introduction}
\label{sec:introduction}

In silico prediction of absorption, distribution, metabolism, excretion, and toxicity (ADMET) properties is now embedded in compound prioritization, where large public task collections and accessible prediction platforms make it possible to screen many endpoints before experimental follow-up \citep{xiong2021admetlab,wei2022interpretable,tian2022admetboost,zhang2022helix,swanson2024admetai,huang2021tdc}. Message-passing networks, geometry-aware models, chemical-language representations, and quantum-informed pretraining have extended the predictive frontier \citep{heid2024chemprop,fang2022geometry,oliveira2022grammar,kim2024quantum}. Yet a better point predictor does not answer a different and operationally important question: when should an individual molecular prediction be treated as informative, uncertain, or in need of deferral?

Aggregate benchmark scores provide an incomplete answer because molecular observations are not independent points in an unstructured feature space. Random splitting can place close analogues on both sides of an evaluation boundary, whereas scaffold and similarity-cluster designs probe different forms of chemical extrapolation. Both predictive performance and uncertainty-method rankings can therefore change with the endpoint, representation, and split mechanism \citep{hirschfeld2020molecularuq,scalia2020gnnuq,lanini2024unique,koetter2025upgrading,tang2025polymeruq,parrondo2026shifts}. Repeated, paired evaluation is also necessary: model families are not independent statistical replicates, and a cross-model average cannot replace uncertainty across matched data partitions \citep{ash2025comparison}.

Reliability itself is multidimensional. Probability calibration asks whether reported probabilities agree with empirical frequencies \citep{guo2017calibration}. Applicability-domain analysis asks whether a query resembles the chemical space represented in training, but molecular novelty, model uncertainty, and realized error are not interchangeable \citep{kim2025distance,koetter2025upgrading}. Selective prediction asks whether rejecting uncertain cases lowers risk, yet in an imbalanced endpoint the same curve can improve because the rejection rule removes one class faster than another. Modern uncertainty reviews consequently distinguish uncertainty sources and require the quality of an uncertainty score to be evaluated rather than inferred from predictive accuracy \citep{hullermeier2021aleatoric,gawlikowski2023survey}. Molecular applications combining uncertainty with attribution or reaction-property prediction further show that useful uncertainty can be local and task specific \citep{komissarov2025explainable,liu2024reactionuq,jiang2024gnas}.

Conformal prediction supplies prediction sets or intervals with finite-sample marginal coverage under exchangeability and has been adapted to toxicity, ADMET, graph models, and chemical covariate shift \citep{zhang2021toxicity,li2024conformalized,laghuvarapu2023codrug}. Marginal validity, however, does not imply class-conditional or subgroup validity. In a rare-event classifier, an overall coverage value close to its nominal target can coexist with near-total failure for the minority class \citep{krstajic2021critical,cauchois2021knowing,bastani2022multivalid,ding2023classconditional}. Conversely, a class-conditional procedure can raise coverage by returning both labels for most compounds. Coverage and informativeness must therefore be reported together \citep{fisch2022trading}.

Molecular distribution shift makes this balance more difficult. Weighted conformal prediction has a clear justification when a target-to-source likelihood ratio is known \citep{tibshirani2019covariate}; in practice, molecular ratios are estimated from finite unlabeled samples in a chosen representation. An estimated ratio can collapse effective sample size, saturate clipping bounds, or reduce one domain discrepancy while worsening another. Adaptive conformal methods address other nonexchangeable settings \citep{gibbs2021adaptive,angelopoulos2023pid}, but algorithmic complexity alone does not establish a useful coverage--efficiency trade-off. What is missing from many empirical workflows is an explicit certificate for when a transport correction is permitted to affect the final set and an exact fallback for when that certificate fails.

Recent conformal methods pursue related but distinct objectives. KMM-CP uses kernel mean matching and support-overlap selection for covariate-shift correction, including molecular benchmarks, while multi-distribution robust conformal prediction assumes labeled data from multiple source distributions and targets uniform validity across those sources or their mixtures \citep{laghuvarapu2026kmmcp,yang2026multidistribution}. Singleton-optimized conformal prediction redesigns the score to increase singleton frequency, whereas batch conformal prediction constructs joint label-vector sets for multiple test examples \citep{wang2025socop,gazin2025batch}. TAME addresses neither multi-source robustness nor joint batch labels: it produces per-compound binary sets for one unlabeled target cohort, admits transport corrections only after view-specific audits, and expands a protected baseline rather than optimizing it away. These differences make the fallback and the measured coverage--deferral trade-off central to its evaluation.

Broad contemporary ADMET benchmarks answer a complementary question by comparing foundation models, graph networks, AutoML systems, and classical predictors under data scarcity, structural out-of-distribution evaluation, class imbalance, activity cliffs, and beyond-Rule-of-5 space \citep{zhao2026revisiting}. Our focus is not another architecture leaderboard. We hold representation and predictor families fixed in the diagnostic stage, then ask whether commonly reported reliability lenses agree and whether their observed failures can be converted into design constraints for an auditable uncertainty layer.

The study has two prospectively separated stages (Figure~\ref{fig:evidence_chain}). Stage~I is a frozen four-endpoint audit of BBBP and ClinTox classification and ESOL and Lipophilicity regression under random, label-blind scaffold, and label-blind similarity-cluster splits. It tests four questions: how molecular distribution shift changes point performance and calibration; when chemical similarity ranks error or miscoverage; whether aggregate coverage conceals class failure; and whether conformal or selective remedies remain informative. Stage~II translates the resulting constraints into the transport-audited multi-view envelope (TAME) and evaluates this method under the separately frozen RACER-C4 protocol. TAME uses an ordinary class-conditional Mondrian set as a protected baseline, two label-free transport views that must pass prespecified audits, and an all-view consensus that may add a protected label but may never remove a baseline label or create a new confident singleton.

Architecture development was restricted to the public Tox21 leaderboard cohort \citep{huang2016tox21challenge}. The method, primary endpoints, five fresh seeds, promotion criteria, a $-5$ percentage-point point-estimate guardrail for Macro correct-singleton yield (MacroCSY), and the negative-result policy were frozen before the independent EPA labels were opened. Final predictions and transport audits were written and hashed before the label file was acquired and parsed. This ordering separates a useful empirical result from a stronger claim that the design cannot support: TAME is evaluated as an audited, failure-aware layer in a locked Tox21 setting, not as an exact conditional-coverage method or a universal solution to arbitrary molecular shift.

\begin{figure}[!htbp]
\centering
\includegraphics[width=\textwidth,height=0.78\textheight,keepaspectratio]{figure_1_evidence_chain.pdf}
\caption{\textbf{Evidence chain from reliability diagnosis to independent validation.} (A) A frozen four-endpoint audit compared point performance, calibration, conformal validity, chemical-domain diagnostics, and selective retention under three label-blind split designs. (B) The audit exposed failure signals that aggregate summaries concealed. (C) These signals became explicit TAME design constraints: two audited label-free transport views, a baseline-containing consensus, and exact fallback. (D) Under the RACER-C4 evaluation protocol, final EPA predictions were hashed and sealed before labels were opened; the primary result was a gain in minimum-class coverage with the MacroCSY point estimate inside the frozen efficiency guardrail.}
\label{fig:evidence_chain}
\end{figure}
\FloatBarrier
